%%%%%%%%%%%%%%%%%%%%%%%%%%%%%%%%%%%%%%%%%%%%%%%%%%%%%%%%%%%%%%%%%%%%%%%%%%%
% Chapter: Analog-to-Digital Conversion
% Falstad examples: examples/adc (generated by falstad/circuits/adc.py)
% Common values (capstone): fs = 5 Hz, 3 bit, full scale 4.00 V, LSB 0.5 V,
% ladder 8 x 1 kOhm from 4 V, Ck = (Vin > k*0.5 V), LEDs with 330 Ohm
%%%%%%%%%%%%%%%%%%%%%%%%%%%%%%%%%%%%%%%%%%%%%%%%%%%%%%%%%%%%%%%%%%%%%%%%%%%

%%%%%%%%%%%%%%%%%%%%%%%%%%%%%%%%%%%%%%%%%%%%%%%%%%%%%%%%%%%%%%%%%%%%%%%%%%%
\section{Sampling and Aliasing}
%%%%%%%%%%%%%%%%%%%%%%%%%%%%%%%%%%%%%%%%%%%%%%%%%%%%%%%%%%%%%%%%%%%%%%%%%%%
\nsl{An analog-to-digital converter (ADC) does two things: it samples the
signal in time, and it quantizes each sample in amplitude. Both steps lose
information, and both losses can be controlled by design. This section
deals with sampling: how fast must we sample, and what happens if the
signal contains frequencies that are too high for the sampling rate? The
answer is the sampling theorem of Nyquist and Shannon \cite{shannon1949}.
Its practical consequence is a filter that every sampled system needs in
front of the S\&H stage: the anti-aliasing low-pass.\newpage}

\begin{description}
\item[{\rot\bf The sampling theorem}]~\\[-\lblskip]
\bi
	\ite a signal is sampled at the instants $t_k = k\,T_s$ with the sampling frequency $f_s = 1/T_s$
\ei
\keyeq{Sampling theorem (Nyquist--Shannon)}{A signal that contains no frequencies at or above
$f_s/2$ is completely determined by its samples taken at the rate $f_s$.
$f_s/2$ is called the \textbf{Nyquist frequency}.}
\bi
	\ite the condition is $f_\mathrm{max} < f_s/2$, i.e.\ more than two samples per period of the fastest component
	\ite it applies to \textbf{every} component of the input, also to noise and interference
	\ite capstone: $f_s = \SI{5}{\hertz}$ $\Rightarrow$ only components below \SI{2.5}{\hertz} are allowed
\ei
\nsl{The theorem says that the original signal can be reconstructed from
its samples, provided that the band limit is respected. If the condition
is violated, the samples are still taken, but they are the samples of a
different, lower frequency. No processing after the sampler can tell the
two cases apart.}
\os{\newpage}

\item[{\rot\bf Aliasing}]~\\[-\lblskip]
\begin{center}
\begin{tikzpicture}
\begin{axis}[width=12cm,height=4.4cm,xmin=0,xmax=2,ymin=-1.3,ymax=1.3,
	xlabel={$t$ / s},axis lines=left,font=\scriptsize,ytick={-1,0,1},
	xtick={0,0.2,...,2.0},xticklabel style={/pgf/number format/fixed},
	legend style={at={(1,1.12)},anchor=south east,draw=none,fill=none},legend columns=3]
\addplot[rwucyan,domain=0:2,samples=400] {cos(deg(2*pi*4*x))};
\addlegendentry{signal \SI{4}{\hertz}}
\addplot[rwuviolet,thick,dashed,domain=0:2,samples=100] {cos(deg(2*pi*x))};
\addlegendentry{alias \SI{1}{\hertz}}
\addplot[only marks,mark=*,mark size=2pt,black,domain=0:2,samples=11] {cos(deg(2*pi*4*x))};
\addlegendentry{samples at \SI{5}{\hertz}}
\end{axis}
\end{tikzpicture}
\end{center}
\bi
	\ite a \SI{4}{\hertz} cosine sampled at \SI{5}{\hertz} gives the same samples as a \SI{1}{\hertz} cosine
\ei
\keyeq{Alias frequency}{$\displaystyle f_\mathrm{alias} = |f - k\,f_s|$, \quad with the integer $k$ that makes $f_\mathrm{alias} \leq f_s/2$}
\os{\newpage}

\item[{\rot\bf Folding}]~\\[-\lblskip]
\begin{center}
\begin{tikzpicture}
\begin{axis}[width=11cm,height=4cm,xmin=0,xmax=21,ymin=0,ymax=3,
	xlabel={input frequency $f$ / Hz},ylabel={$f_\mathrm{alias}$ / Hz},axis lines=left,font=\scriptsize,
	xtick={0,2.5,5,7.5,10,12.5,15,17.5,20},ytick={0,2.5}]
\addplot[rwuviolet,thick] coordinates {(0,0) (2.5,2.5) (5,0) (7.5,2.5) (10,0) (12.5,2.5) (15,0) (17.5,2.5) (20,0) (21,1)};
\fill[rwucyan!25] (axis cs:0,0) rectangle (axis cs:2.5,2.9);
\node[font=\tiny,anchor=north] at (axis cs:1.25,2.9) {pass band};
\end{axis}
\end{tikzpicture}
\end{center}
\bi
	\ite every frequency is folded into the range $0 \ldots f_s/2$
	\ite $f_s = \SI{5}{\hertz}$: \SI{4}{\hertz} $\to$ \SI{1}{\hertz}; \SI{12}{\hertz} $\to$ \SI{2}{\hertz}; \SI{49}{\hertz} $\to$ \SI{1}{\hertz}; \SI{50}{\hertz} $\to$ \SI{0}{\hertz}
	\ite after sampling an alias cannot be distinguished from a real signal component
\ei
\os{\newpage}

\item[{\rot\bf The 50 Hz problem}]~\\[-\lblskip]
\bi
	\ite mains interference (\SI{50}{\hertz}) is coupled into every sensor cable
	\ite capstone: $f_s = \SI{5}{\hertz}$ $\Rightarrow$ $f_\mathrm{alias} = |\SI{50}{\hertz} - 10\cdot\SI{5}{\hertz}| = \SI{0}{\hertz}$
	\bi
		\itee every sample sees the 50 Hz sine at the same phase
		\itee the result is a \textbf{constant DC offset} that biases every reading
		\itee its size depends on the (random) phase: anywhere between $-\hat V$ and $+\hat V$
		\itee nothing on the display looks wrong: the error is invisible
	\ei
	\ite with $f_s = \SI{4.9}{\hertz}$ the alias is $|\SI{50}{\hertz}-10\cdot\SI{4.9}{\hertz}| = \SI{1}{\hertz}$: a slow beat
	\ite \textbf{the anti-aliasing low-pass in front of the S\&H stage is mandatory}
\ei
\nsl{Because $f_s$ is an integer fraction of the mains frequency, the
interference is sampled at the same point of its period every time. The
digital system sees a perfectly stable value, shifted by an unknown
amount. With a slightly different sampling rate the same interference
would show up as a slow oscillation of the readings, which is at least
visible. The only remedy is to remove the interference before sampling.}
\os{\newpage}

\item[{\rot\bf Falstad: sampling 50 Hz hum}]~\\[-\lblskip]
\bi
	\ite input: \SI{2}{\volt} DC (the wanted signal) plus \SI{0.5}{\volt} at \SI{50}{\hertz}
	\ite two S\&H stages ($R = \SI{1}{\kilo\ohm}$, $C_H = \SI{100}{\nano\farad}$, track window \SI{1}{\milli\second})
	\bi
		\itee $f_s = \SI{5}{\hertz}$: every sample is \SI{2.14}{\volt}: an offset of \SI{0.14}{\volt}
		\itee $f_s = \SI{4.9}{\hertz}$: \SI{2.14}{\volt}, \SI{2.50}{\volt}, \SI{2.14}{\volt}, \SI{1.58}{\volt}, \ldots: a \SI{1}{\hertz} sine
	\ei
	\ite change the phase of the \SI{50}{\hertz} source: the offset at \SI{5}{\hertz} changes between \SI{1.5}{\volt} and \SI{2.5}{\volt}
\ei
\falstad{alias50}
\nsl{The track window is only \SI{1}{\milli\second} here, so that the hold
capacitor can follow the \SI{50}{\hertz} signal ($\tau =
\SI{0.1}{\milli\second}$). The held value is the input at the end of the
window, $t = \SI{1}{\milli\second}$:
$\SI{2}{\volt} + \SI{0.5}{\volt}\cdot\sin(2\pi\cdot\SI{50}{\hertz}\cdot
\SI{1}{\milli\second} - \SI{1.8}{\degree}) = \SI{2.14}{\volt}$, where
\SI{1.8}{\degree} is the phase lag of the RC charging path.}
\os{\newpage}

\item[{\rot\bf Anti-aliasing filter}]~\\[-\lblskip]
\begin{center}
\begin{tikzpicture}
\begin{axis}[width=11cm,height=4.2cm,xmode=log,xmin=0.1,xmax=100,ymin=-60,ymax=5,
	xlabel={$f$ / Hz},ylabel={$|H|$ / dB},axis lines=left,font=\scriptsize,grid=major,
	legend style={at={(0.02,0.05)},anchor=south west,font=\tiny}]
\addplot[rwucyan,thick,domain=0.1:100,samples=100] {-10*log10(1+(x/2.34)^2)};
\addlegendentry{1st order, $f_c=\SI{2.34}{\hertz}$}
\addplot[rwuviolet,thick,domain=0.1:100,samples=100] {-10*log10(1+(x/2.34)^4)};
\addlegendentry{2nd order Butterworth}
\addplot[black,dashed] coordinates {(2.5,-60) (2.5,5)};
\addplot[black,dashed] coordinates {(50,-60) (50,5)};
\node[font=\tiny,anchor=west] at (axis cs:2.6,-50) {$f_s/2$};
\node[font=\tiny,anchor=east] at (axis cs:48,-50) {\SI{50}{\hertz}};
\end{axis}
\end{tikzpicture}
\end{center}
\bi
	\ite pass the wanted signal, attenuate everything above $f_s/2$ below the resolution
	\ite at \SI{50}{\hertz}: 1st order \SI{-26.6}{\dB} (factor 21), 2nd order \SI{-53}{\dB} (factor 460)
	\ite the filter must be \textbf{before} the S\&H: after sampling the damage is done
\ei
\os{\newpage}
\end{description}

\designcard{Nyquist \& Anti-aliasing}{%
\begin{itemize}\itemsep 0pt
\item sampling theorem: all input components must be below $f_s/2$ (Nyquist frequency)
\item alias: $f_\mathrm{alias} = |f - k f_s|$, nearest multiple $k f_s$; result lies in $0\ldots f_s/2$
\item capstone: $f_s = \SI{5}{\hertz}$, \SI{50}{\hertz} hum $\to$ $|50 - 10\cdot 5|\,\si{\hertz} = \SI{0}{\hertz}$: constant offset on every reading
\item low-pass before the S\&H: $f_c \approx f_s/2$ (capstone \SI{2.5}{\hertz}), e.g.\ RC: $f_c = 1/(2\pi RC)$
\item required attenuation: interference amplitude $\times |H(f)| < \frac12\,\LSB$ (better $\frac14\,\LSB$)
\item 1st order: $|H| \approx f_c/f$ (\SI{-20}{\dB}/decade), 2nd order: $(f_c/f)^2$ (\SI{-40}{\dB}/decade)
\item higher $f_s$ (oversampling) relaxes the filter: more distance between signal band and $f_s/2$
\end{itemize}}
\os{\newpage}

%%%%%%%%%%%%%%%%%%%%%%%%%%%%%%%%%%%%%%%%%%%%%%%%%%%%%%%%%%%%%%%%%%%%%%%%%%%
\section{Quantization}
%%%%%%%%%%%%%%%%%%%%%%%%%%%%%%%%%%%%%%%%%%%%%%%%%%%%%%%%%%%%%%%%%%%%%%%%%%%
\nsl{The second step is quantization: the continuous voltage range is
divided into $2^N$ intervals of equal width, and each sample is replaced by
the number of its interval. The width of the interval is the smallest
voltage step the converter can resolve; it is called LSB (least
significant bit), because a change of the input by one LSB changes the
least significant bit of the output code \cite{kester2005,tietze2008}.\newpage}

\begin{description}
\item[{\rot\bf Resolution and LSB}]~\\[-\lblskip]
\keyeq{LSB}{$\displaystyle \LSB = \frac{V_\mathrm{FS}}{2^N}$ \qquad ($V_\mathrm{FS}$: full-scale range, $N$: number of bits)}
\bi
	\ite capstone: $N = 3$, $V_\mathrm{FS} = \SI{4}{\volt}$: $\LSB = \SI{0.5}{\volt}$, 8 codes 000 \ldots 111
	\ite 10 bit, \SI{3.3}{\volt} (microcontroller): \SI{3.2}{\milli\volt}; 16 bit, \SI{5}{\volt}: \SI{76}{\micro\volt}
	\ite code for an input voltage (our flash converter):\\
	$\displaystyle D = \left\lfloor \frac{\Vin}{\LSB}\right\rfloor$, \quad limited to $0 \ldots 2^N - 1$
	\ite back to a voltage: $V_D = D\cdot\LSB$ (lower edge) or $(D+\frac12)\,\LSB$ (bin centre)
\ei
\os{\newpage}

\item[{\rot\bf Transfer characteristic}]~\\[-\lblskip]
\begin{center}
\begin{tikzpicture}
\begin{axis}[width=8.4cm,height=5.4cm,xmin=0,xmax=4.2,ymin=0,ymax=7.6,
	xlabel={$\Vin$ / V},ylabel={code $D$},axis lines=left,font=\scriptsize,
	xtick={0,0.5,...,4},ytick={0,...,7},
	yticklabels={000,001,010,011,100,101,110,111},grid=major,grid style={gray!20}]
\addplot[rwuviolet,very thick,const plot] coordinates {(0,0) (0.5,1) (1,2) (1.5,3) (2,4) (2.5,5) (3,6) (3.5,7) (4.2,7)};
\addplot[rwucyan,dashed,domain=0:4] {x/0.5};
\addplot[only marks,mark=o,rwucyandark] coordinates {(0.25,0) (0.75,1) (1.25,2) (1.75,3) (2.25,4) (2.75,5) (3.25,6) (3.75,7)};
\end{axis}
\end{tikzpicture}
\end{center}
\bi
	\ite staircase: thresholds at $k\cdot\LSB$ ($k = 1\ldots 7$), the code changes at \SI{0.5}{\volt}, \SI{1.0}{\volt}, \ldots, \SI{3.5}{\volt}
	\ite circles: bin centres \SI{0.25}{\volt}, \SI{0.75}{\volt}, \ldots, \SI{3.75}{\volt}: safe test points
\ei
\os{\newpage}

\item[{\rot\bf Quantization error}]~\\[-\lblskip]
\begin{center}
\begin{tikzpicture}
\begin{axis}[width=8.4cm,height=3.6cm,xmin=0,xmax=4.2,ymin=-0.1,ymax=0.6,
	xlabel={$\Vin$ / V},ylabel={$\Vin - D\cdot\LSB$ / V},axis lines=left,font=\scriptsize,
	xtick={0,0.5,...,4},ytick={0,0.25,0.5}]
\addplot[rwuviolet,thick] coordinates {(0,0) (0.5,0.5) (0.5,0) (1,0.5) (1,0) (1.5,0.5) (1.5,0) (2,0.5) (2,0) (2.5,0.5) (2.5,0) (3,0.5) (3,0) (3.5,0.5) (3.5,0) (4,0.5)};
\end{axis}
\end{tikzpicture}
\end{center}
\bi
	\ite with thresholds at $k\cdot\LSB$: error $0 \ldots 1\,\LSB$ (truncation), mean $\frac12\,\LSB$
	\ite shift all thresholds by $\frac12\,\LSB$ (ladder: $R/2$ at the bottom, $3R/2$ at the top): error $\pm\frac12\,\LSB$
	\ite or read the code as the bin centre $(D+\frac12)\,\LSB$: also $\pm\frac12\,\LSB$
\ei
\nsl{The two variants differ only by a constant offset of half an LSB.
The capstone uses the simple ladder of eight equal resistors; its
thresholds are at integer multiples of \SI{0.5}{\volt}. A digital value
$D$ then means ``the input is between $D\cdot\SI{0.5}{\volt}$ and
$(D+1)\cdot\SI{0.5}{\volt}$''.}
\os{\newpage}

\item[{\rot\bf Quantization noise and SNR}]~\\[-\lblskip]
\bi
	\ite for a busy signal the error behaves like noise, uniformly distributed over one LSB
	\bi
		\itee rms value $\displaystyle V_{q,\mathrm{rms}} = \frac{\LSB}{\sqrt{12}}$
	\ei
	\ite full-scale sine: $V_\mathrm{rms} = \frac{V_\mathrm{FS}}{2\sqrt 2} = \frac{2^N\,\LSB}{2\sqrt2}$
\ei
\keyeq{Signal-to-noise ratio of an ideal $N$-bit converter}{$\displaystyle \mathrm{SNR} = 20\log_{10}\frac{2^N\sqrt{12}}{2\sqrt 2}\,\mathrm{dB} = 6.02\,N + 1.76\dB$}
\bi
	\ite 3 bit: \SI{19.8}{\dB}; 8 bit: \SI{49.9}{\dB}; 12 bit: \SI{74.0}{\dB}; 16 bit: \SI{98.1}{\dB}
	\ite every additional bit: \SI{6}{\dB} more (factor 2 in amplitude)
\ei
\os{\newpage}

\item[{\rot\bf Falstad: quantization made visible}]~\\[-\lblskip]
\bi
	\ite 3-bit flash converter (next sections) followed by a 3-bit R-2R DAC
	\bi
		\itee the DAC turns the code back into $V_q = D\cdot\SI{0.5}{\volt}$
	\ei
	\ite input: triangle \SIrange{0}{4}{\volt}, \SI{0.25}{\hertz}
	\ite the scope shows $\Vin$ and the staircase $V_q$; the difference is always between \SI{0}{\volt} and \SI{0.5}{\volt}
	\ite example readings: $\Vin = \SI{1.8}{\volt}$ $\to$ $V_q = \SI{1.5}{\volt}$; $\Vin = \SI{3.98}{\volt}$ $\to$ $V_q = \SI{3.5}{\volt}$
\ei
\falstad{quant_error}
\os{\newpage}
\end{description}

\designcard{Reference Ladder \& Quantization}{%
\begin{itemize}\itemsep 0pt
\item $\LSB = V_\mathrm{FS}/2^N$; capstone: $\SI{4}{\volt}/8 = \SI{0.5}{\volt}$
\item ladder of $2^N$ equal resistors from $\Vref = V_\mathrm{FS}$ to ground: taps $V_k = k\cdot\Vref/2^N$, $k = 1\ldots 2^N-1$
\item capstone: $8\times\SI{1}{\kilo\ohm}$ from \SI{4}{\volt}: taps \SI{0.5}{\volt}, \SI{1.0}{\volt}, \ldots, \SI{3.5}{\volt}; $I = \SI{0.5}{\milli\ampere}$
\item code $D = \lfloor\Vin/\LSB\rfloor$; quantization error $0\ldots 1\,\LSB$; with $R/2$ at the bottom and $3R/2$ at the top: $\pm\frac12\,\LSB$
\item test points: bin centres $(D+\frac12)\,\LSB$ = \SI{0.25}{\volt}, \SI{0.75}{\volt}, \ldots, \SI{3.75}{\volt}
\item ladder current $\gg$ comparator input currents; tap resistance $\leq R_\mathrm{total}/4 = \SI{2}{\kilo\ohm}$
\item reference: stiff voltage source (reference IC + buffer); its error scales all thresholds
\item SNR of an ideal converter: $6.02\,N + 1.76\dB$
\end{itemize}}
\os{\newpage}

%%%%%%%%%%%%%%%%%%%%%%%%%%%%%%%%%%%%%%%%%%%%%%%%%%%%%%%%%%%%%%%%%%%%%%%%%%%
\section{The Comparator}
%%%%%%%%%%%%%%%%%%%%%%%%%%%%%%%%%%%%%%%%%%%%%%%%%%%%%%%%%%%%%%%%%%%%%%%%%%%
\nsl{The comparator is the 1-bit converter: it decides whether its input is
above or below a threshold. Every converter contains at least one. An
op-amp without feedback is a comparator, because its huge open-loop gain
drives the output into one of the rails as soon as the input difference is
more than a few microvolts. Dedicated comparator ICs are faster and have
logic-compatible outputs \cite{horowitz2015,tietze2008}.\newpage}

\begin{description}
\item[{\rot\bf Open-loop op-amp as comparator}]~\\[-\lblskip]
\begin{center}
\begin{circuitikz}[scale=0.85, transform shape]
  \draw (0,0) node[op amp, noinv input up](k){};
  \draw (k.+) -- ++(-1,0) node[left]{$\Vin$};
  \draw (k.-) -- ++(-1,0) node[left]{$\Vref$};
  \draw (k.out) to[short,-o] ++(0.8,0) node[right]{$\Vout$: \SI{0}{\volt} / \SI{5}{\volt}};
  \draw (k.up) -- ++(0,0.4) node[vcc]{\SI{5}{\volt}};
  \draw (k.down) -- ++(0,-0.3) node[ground]{};
\end{circuitikz}
\begin{tikzpicture}
\begin{axis}[width=5cm,height=3.4cm,xmin=0,xmax=4,ymin=-0.5,ymax=5.8,
	xlabel={$\Vin$ / V},ylabel={$\Vout$ / V},axis lines=left,font=\scriptsize,xtick={0,2,4},ytick={0,5}]
\addplot[rwuviolet,very thick] coordinates {(0,0) (2,0) (2,5) (4,5)};
\node[font=\tiny,anchor=west] at (axis cs:2.05,2.5) {$\Vref=\SI{2}{\volt}$};
\end{axis}
\end{tikzpicture}
\end{center}
\bi
	\ite $\Vout = V_\mathrm{OH}$ if $\Vin > \Vref$, otherwise $V_\mathrm{OL}$
	\ite single \SI{5}{\volt} supply: output levels fit 5 V logic directly (rail-to-rail output)
	\ite Falstad: op-amp with \emph{max output} \SI{5}{\volt}, \emph{min output} \SI{0}{\volt}
\ei
\os{\newpage}

\item[{\rot\bf Real comparators}]~\\[-\lblskip]
\bi
	\ite op-amps are slow when used open loop: recovery from saturation, slew rate (µs)
	\ite comparator ICs (LM393, LM339, TLV3501): fast, no frequency compensation, logic outputs
	\bi
		\itee LM393/LM339: \textbf{open-collector} output, needs a pull-up resistor (e.g.\ \SI{4.7}{\kilo\ohm} to \SI{5}{\volt})
	\ei
	\ite input offset voltage (mV): shifts the threshold; must be $\ll \LSB$
	\ite input bias current: flows through the ladder resistors, small threshold shift
	\ite never use a comparator IC as an op-amp: it has no compensation and oscillates with feedback
\ei
\nsl{For the 3-bit capstone with an LSB of \SI{0.5}{\volt}, any op-amp or
comparator is good enough: offsets of a few millivolts are three orders of
magnitude below the LSB. For an 8-bit flash converter with
\SI{4}{\volt} full scale ($\LSB = \SI{15.6}{\milli\volt}$) the offset of
each comparator already matters.}
\os{\newpage}

\item[{\rot\bf Chatter at slow, noisy inputs}]~\\[-\lblskip]
\bi
	\ite a slowly changing input with noise crosses the threshold several times
	\ite the output switches back and forth (\textbf{chatter}): a counter counts wrong, an LED flickers
	\ite remedy: \textbf{hysteresis}: two thresholds, $V_{T+}$ for rising and $V_{T-}$ for falling inputs
	\bi
		\itee positive feedback (Schmitt trigger)
		\itee the hysteresis $V_{T+} - V_{T-}$ must be larger than the peak-to-peak noise
	\ei
	\ite Falstad: triangle \SIrange{0}{4}{\volt} at \SI{1}{\hertz} plus \SI{0.2}{\volt} at \SI{50}{\hertz}
	\bi
		\itee plain comparator: 10 edges in one period; Schmitt trigger: exactly 2
	\ei
\ei
\falstad{comparator_hyst}
\os{\newpage}

\item[{\rot\bf Schmitt trigger (non-inverting)}]~\\[-\lblskip]
\begin{center}
\begin{circuitikz}[scale=0.8, transform shape]
  \draw (0,0) node[op amp, noinv input up](k){};
  \draw (k.+) -- ++(-0.5,0) coordinate(p) to[R, l_=$R_1$] ++(-2,0) node[left]{$\Vin$};
  \draw (p) to[short,*-] ++(0,1.2) coordinate(q) to[R, l=$R_2$] (q -| k.out) -- ++(0.3,0) |- (k.out);
  \draw (k.-) -- ++(-0.8,0) node[left]{$\Vref$};
  \draw (k.out) ++(0.3,0) node[circ]{} to[short,-o] ++(0.8,0) node[right]{$\Vout$};
\end{circuitikz}
\begin{tikzpicture}
\begin{axis}[width=5cm,height=3.4cm,xmin=0,xmax=4,ymin=-0.5,ymax=5.8,
	xlabel={$\Vin$ / V},ylabel={$\Vout$ / V},axis lines=left,font=\scriptsize,xtick={0,1.7,2.2,4},
	xticklabels={0,1.7,2.2,4},ytick={0,5}]
\addplot[rwuviolet,very thick,-{Stealth}] coordinates {(0.3,0) (2.2,0) (2.2,5) (3.8,5)};
\addplot[rwucyan,very thick,-{Stealth}] coordinates {(3.8,5.15) (1.7,5.15) (1.7,0.15) (0.3,0.15)};
\end{axis}
\end{tikzpicture}
\end{center}
\bi
	\ite switching when $V_+ = \Vref$, with $V_+ = \Vin\frac{R_2}{R_1+R_2} + \Vout\frac{R_1}{R_1+R_2}$:\\
	$\displaystyle V_{T} = \Vref\left(1+\frac{R_1}{R_2}\right) - \Vout\,\frac{R_1}{R_2}$
	\ite $V_{T+}$ (output low): $\Vref(1+R_1/R_2)$; \quad $V_{T-}$ (output high): $V_{T+} - V_\mathrm{OH}R_1/R_2$
	\ite $R_1 = \SI{10}{\kilo\ohm}$, $R_2 = \SI{100}{\kilo\ohm}$, $\Vref = \SI{2}{\volt}$: $V_{T+} = \SI{2.2}{\volt}$, $V_{T-} = \SI{1.7}{\volt}$, hysteresis \SI{0.5}{\volt}
\ei
\os{\newpage}
\end{description}

\designcard{Comparator}{%
\begin{itemize}\itemsep 0pt
\item $\Vout = $ high if $V_+ > V_-$, else low; op-amp open loop or comparator IC
\item logic compatible: supply $0/\SI{5}{\volt}$, rail-to-rail output (or open collector + pull-up \SI{4.7}{\kilo\ohm})
\item Falstad: op-amp, max output \SI{5}{\volt}, min output \SI{0}{\volt}; logic inputs switch at \SI{2.5}{\volt}
\item offset voltage $\ll \LSB$; input current $\times$ source resistance $\ll \LSB$
\item hysteresis against chatter: $R_1$ from $\Vin$ to $V_+$, $R_2$ from $\Vout$ to $V_+$, $\Vref$ at $V_-$
\item $V_{T+} = \Vref(1+R_1/R_2)$, $V_{T-} = V_{T+} - V_\mathrm{OH}R_1/R_2$, hysteresis $V_\mathrm{OH}R_1/R_2 >$ noise (peak-to-peak)
\item in a flash ADC behind an S\&H: no hysteresis needed (the held input is constant), hysteresis would shift the thresholds
\end{itemize}}
\os{\newpage}

%%%%%%%%%%%%%%%%%%%%%%%%%%%%%%%%%%%%%%%%%%%%%%%%%%%%%%%%%%%%%%%%%%%%%%%%%%%
\section{The Flash Converter}
%%%%%%%%%%%%%%%%%%%%%%%%%%%%%%%%%%%%%%%%%%%%%%%%%%%%%%%%%%%%%%%%%%%%%%%%%%%
\nsl{The flash (or parallel) converter is the most direct way to convert a
voltage: compare it with all thresholds at the same time. An $N$-bit flash
converter needs $2^N - 1$ comparators and a reference ladder that provides
the thresholds. The comparator outputs form a thermometer code, which a
small logic circuit turns into a binary number. The conversion takes only
the delay of one comparator and a few gates, which makes flash converters
the fastest of all architectures. The price is the exponential growth of
the number of comparators: 7 for 3 bit, 255 for 8 bit.\newpage}

\begin{description}
\item[{\rot\bf 2-bit flash converter}]~\\[-\lblskip]
\begin{center}
\begin{circuitikz}[scale=0.62, transform shape]
  \draw (0,8) node[above]{$\Vref = \SI{4}{\volt}$} node[circ]{}
        to[R, l_=$R$] (0,6) to[R, l_=$R$] (0,4) to[R, l_=$R$] (0,2) to[R, l_=$R$] (0,0) node[ground]{};
  \foreach \k/\y in {3/6,2/4,1/2} {
    \draw (3.2,\y) node[op amp, noinv input down, anchor=-, scale=0.8](c\k){};
    \draw (0,\y) node[circ]{} node[left=6mm]{\SI{\k}{\volt}} -- (c\k.-);
    \draw (c\k.+) -- ++(-0.4,0) node[left]{\small $\Vin$};
    \draw (c\k.out) -- ++(0.4,0) node[right]{$C_\k$};
  }
  \node[draw,thick,minimum width=2.6cm,minimum height=4.4cm,align=center] (enc) at (9.2,3.5) {encoder\\[2mm]$B_1 = C_2$\\[1mm]$B_0 = C_1\overline{C_2} + C_3$};
  \draw[->] (enc.east) ++(0,0.5) -- ++(0.8,0) node[right]{$B_1$};
  \draw[->] (enc.east) ++(0,-0.5) -- ++(0.8,0) node[right]{$B_0$};
\end{circuitikz}
\end{center}
\bi
	\ite $\LSB = \SI{1}{\volt}$; thresholds \SI{1}{\volt}, \SI{2}{\volt}, \SI{3}{\volt}; 3 comparators
	\ite $C_k = 1$ if $\Vin > k\cdot\LSB$
\ei
\falstad{flash2}
\os{\newpage}

\item[{\rot\bf 3-bit flash converter (capstone)}]~\\[-\lblskip]
\begin{center}
\begin{circuitikz}[x=0.62cm,y=0.62cm,font=\scriptsize]
  \ctikzset{bipoles/length=0.5cm}
  \draw (0,9.2) node[above]{$\Vref = \SI{4}{\volt}$} node[circ]{} -- (0,8.6)
        to[R] (0,7.6) to[R] (0,6.6) to[R] (0,5.6) to[R] (0,4.6) to[R] (0,3.6)
        to[R] (0,2.6) to[R] (0,1.6) to[R] (0,0.6) -- (0,0.3) node[ground]{};
  \foreach \k in {1,...,7} {
    \pgfmathsetmacro{\v}{0.5*\k}
    \pgfmathsetmacro{\y}{0.6+\k}
    \draw (3.4,\y) node[op amp, noinv input down, anchor=-, scale=0.42](c\k){};
    \draw (0,\y) node[circ]{} -- (c\k.-);
    \draw (c\k.+) -- ++(-0.4,0) node[left]{$\Vin$};
    \draw (c\k.out) -- ++(0.4,0) node[right]{$C_{\k}$};
    \node[left] at (-0.1,\y) {\pgfmathprintnumber[fixed,precision=1,fixed zerofill]{\v}\,V};
  }
  \node[anchor=west] at (0.4,0.0) {$8\times R$, $R = \SI{1}{\kilo\ohm}$};
  \node[draw,thick,minimum width=1.8cm,minimum height=4.6cm,align=center] (enc) at (9.6,4.6) {thermo-\\meter\\to binary\\encoder};
  \foreach \k in {1,...,7} { \draw[->] (6.6,0.6+\k) -- (enc.west |- 0,0.6+\k); }
  \draw[->] (enc.east |- 0,5.6) -- ++(0.8,0) node[right]{$B_2$};
  \draw[->] (enc.east |- 0,4.6) -- ++(0.8,0) node[right]{$B_1$};
  \draw[->] (enc.east |- 0,3.6) -- ++(0.8,0) node[right]{$B_0$};
\end{circuitikz}
\end{center}
\bi
	\ite 7 comparators, $C_k = (\Vin > k\cdot\SI{0.5}{\volt})$; ladder current $\SI{4}{\volt}/\SI{8}{\kilo\ohm} = \SI{0.5}{\milli\ampere}$
\ei
\os{\newpage}

\item[{\rot\bf Thermometer code}]~\\[-\lblskip]
\begin{center}
\small
\begin{tabular}{cc|ccccccc|ccc}
\toprule
$\Vin$ / V & $D$ & $C_7$ & $C_6$ & $C_5$ & $C_4$ & $C_3$ & $C_2$ & $C_1$ & $B_2$ & $B_1$ & $B_0$\\
\midrule
0.00 \ldots 0.50 & 0 & 0&0&0&0&0&0&0 & 0&0&0\\
0.50 \ldots 1.00 & 1 & 0&0&0&0&0&0&1 & 0&0&1\\
1.00 \ldots 1.50 & 2 & 0&0&0&0&0&1&1 & 0&1&0\\
1.50 \ldots 2.00 & 3 & 0&0&0&0&1&1&1 & 0&1&1\\
2.00 \ldots 2.50 & 4 & 0&0&0&1&1&1&1 & 1&0&0\\
2.50 \ldots 3.00 & 5 & 0&0&1&1&1&1&1 & 1&0&1\\
3.00 \ldots 3.50 & 6 & 0&1&1&1&1&1&1 & 1&1&0\\
3.50 \ldots 4.00 & 7 & 1&1&1&1&1&1&1 & 1&1&1\\
\bottomrule
\end{tabular}
\end{center}
\bi
	\ite like the column of a thermometer: all comparators below the input level are 1
	\ite the number of ones is the code $D$; only $2^N$ of the $2^{2^N-1}$ patterns are valid
\ei
\os{\newpage}

\item[{\rot\bf Deriving the encoder}]~\\[-\lblskip]
\bi
	\ite the top of the column: $H_k = C_k\,\overline{C_{k+1}}$ is 1 only for $D = k$ (one-hot code, $C_8 := 0$)
	\ite each output bit is the OR of the $H_k$ whose code has this bit set:
	\bi
		\itee $B_2 = H_4 + H_5 + H_6 + H_7$ \hfill $= C_4$
		\itee $B_1 = H_2 + H_3 + H_6 + H_7 = C_2\overline{C_3} + C_3\overline{C_4} + C_6\overline{C_7} + C_7$ \hfill $= C_2\overline{C_4} + C_6$
		\itee $B_0 = H_1 + H_3 + H_5 + H_7$ \hfill $= C_1\overline{C_2} + C_3\overline{C_4} + C_5\overline{C_6} + C_7$
	\ei
	\ite simplification uses the thermometer property: $C_{k+1} = 1 \Rightarrow C_k = 1$
	\bi
		\itee $C_2\overline{C_3} + C_3\overline{C_4} = C_2\overline{C_4}$: ``between level 2 and level 4''
		\itee $B_2$: four ones in a row means $D \geq 4$
	\ei
\ei
\nsl{The derivation can also be read directly from the table: $B_2$ is 1
for the lower half of the table, exactly where $C_4$ is 1. $B_1$ is 1 for
the codes 2, 3 and 6, 7: the input is above level 2 but not above level 4,
or it is above level 6. $B_0$ alternates with every code; it is 1 if the
top of the thermometer column is at an odd position. The equations need
three inverters, four 2-input AND gates, one 2-input OR gate and one
4-input OR gate.}
\os{\newpage}

\item[{\rot\bf Encoder logic}]~\\[-\lblskip]
\begin{center}
\begin{circuitikz}[scale=0.72, transform shape]
  % B2
  \draw (0,8.4) node[left]{$C_4$} -- (9.5,8.4) node[right]{$B_2$};
  % B1
  \node[american not port] (n4) at (1.6,6.4) {};
  \draw (n4.in) -- ++(-0.4,0) node[left]{$C_4$};
  \node[american and port] (a24) at (4.6,7.1) {};
  \draw (a24.in 1) -- (0,7.1 |- a24.in 1) node[left]{$C_2$};
  \draw (n4.out) -- ++(0.4,0) |- (a24.in 2);
  \node[american or port] (o1) at (8.2,6.8) {};
  \draw (a24.out) -- ++(0.4,0) |- (o1.in 1);
  \draw (o1.in 2) -- ++(-0.4,0) |- (5,5.9) node[left]{$C_6$};
  \draw (o1.out) -- (9.5,6.8) node[right]{$B_1$};
  % B0
  \foreach \i/\a/\b/\y in {1/1/2/5.1, 3/3/4/3.5, 5/5/6/1.9} {
    \node[american not port] (m\i) at (1.6,\y-0.7) {};
    \draw (m\i.in) -- ++(-0.4,0) node[left]{$C_\b$};
    \node[american and port] (b\i) at (4.6,\y) {};
    \draw (b\i.in 1) -- (0,\y |- b\i.in 1) node[left]{$C_\a$};
    \draw (m\i.out) -- ++(0.4,0) |- (b\i.in 2);
  }
  \node[american or port, number inputs=4] (o0) at (8.2,3.3) {};
  \draw (b1.out) -- ++(0.5,0) |- (o0.in 1);
  \draw (b3.out) -- ++(0.8,0) |- (o0.in 2);
  \draw (b5.out) -- ++(0.5,0) |- (o0.in 3);
  \draw (o0.in 4) -- ++(-0.3,0) |- (5,0.7) node[left]{$C_7$};
  \draw (o0.out) -- (9.5,3.3) node[right]{$B_0$};
\end{circuitikz}
\end{center}
\bi
	\ite 3 NOT, 4 AND, 2 OR gates: the same gates as in the Falstad model
\ei
\os{\newpage}

\item[{\rot\bf Falstad: 3-bit flash converter with display}]~\\[-\lblskip]
\bi
	\ite input: slider \SIrange{0}{4}{\volt} (label \texttt{Vin}); ladder $8\times\SI{1}{\kilo\ohm}$, \SI{4}{\volt}
	\ite comparator outputs and encoder signals are labeled nodes \texttt{C1}\ldots\texttt{C7}, \texttt{nC2}, \texttt{a24}, \ldots, \texttt{B0}\ldots\texttt{B2}
	\ite 7 LEDs as bar graph, 3 LEDs for $B_2B_1B_0$, each with \SI{330}{\ohm}
	\ite checked at the bin centres \SI{0.25}{\volt}, \SI{0.75}{\volt}, \ldots, \SI{3.75}{\volt}: codes 000, 001, \ldots, 111
	\ite default \SI{2.6}{\volt}: five bar LEDs on, binary 101
\ei
\falstad{flash3}
\nsl{Labeled nodes keep the drawing readable: every comparator output gets
a label, and the gate inputs carry the same label names instead of long
wires. Hover over a label to see its voltage. The tap voltages of the
ladder are \SI{0.5}{\volt} \ldots \SI{3.5}{\volt}, the ladder current is
\SI{0.50}{\milli\ampere}.}
\os{\newpage}

\item[{\rot\bf Priority encoder and bubble errors}]~\\[-\lblskip]
\bi
	\ite \textbf{bubble}: a wrong comparator decision inside the column, e.g.\ by offset or noise near a threshold
	\bi
		\itee true code 3: 0000111; with bubble at $C_2$: 0000101
		\itee our encoder: $B_0 = C_1\overline{C_2} = 1$, $B_1 = 0$, $B_2 = 0$ $\Rightarrow$ 001: error of 2 LSB
	\ei
	\ite \textbf{priority encoder} (74HC148, 8 inputs $\to$ 3 bit): outputs the number of the \textbf{highest} active input
	\bi
		\itee ignores everything below the top: a bubble does no harm
	\ei
	\ite bubble correction with 3-input gates: $H_k = C_{k-1}\,C_k\,\overline{C_{k+1}}$
	\ite large flash converters: ones counter (Wallace tree) or Gray-code encoding
\ei
\nsl{In the capstone the comparators decide on a held, constant voltage
with an LSB of \SI{0.5}{\volt}, so bubbles do not occur. In fast converters
with hundreds of comparators they are a real problem, because the
comparators near the input level see differences of only microvolts and
decide at slightly different times.}
\os{\newpage}
\end{description}

\designcard{Thermometer-to-Binary Encoder}{%
\begin{itemize}\itemsep 0pt
\item thermometer code: $C_k = 1$ for all $k \leq D$; the number of ones is the code $D$
\item one-hot: $H_k = C_k\overline{C_{k+1}}$ (top of the column); $B_i = $ OR of all $H_k$ with bit $i$ of $k$ set
\item 3 bit (capstone):\quad $B_2 = C_4$\quad $B_1 = C_2\overline{C_4} + C_6$\quad
	$B_0 = C_1\overline{C_2} + C_3\overline{C_4} + C_5\overline{C_6} + C_7$
\item gates: 3 NOT ($\overline{C_2}, \overline{C_4}, \overline{C_6}$), 4 AND2, 1 OR2, 1 OR4
\item check: sweep $\Vin$ over the bin centres \SIrange{0.25}{3.75}{\volt} in steps of \SI{0.5}{\volt}: 000 \ldots 111
\item alternative: priority encoder (74HC148, active-low inputs and outputs), robust against bubbles
\end{itemize}}
\os{\newpage}

%%%%%%%%%%%%%%%%%%%%%%%%%%%%%%%%%%%%%%%%%%%%%%%%%%%%%%%%%%%%%%%%%%%%%%%%%%%
\section{Displaying the Result}
%%%%%%%%%%%%%%%%%%%%%%%%%%%%%%%%%%%%%%%%%%%%%%%%%%%%%%%%%%%%%%%%%%%%%%%%%%%
\nsl{In the lab the digital value is shown on LEDs: a bar graph driven
directly by the thermometer code, and three LEDs for the binary code. The
bar graph is easy to read without any decoding; the binary LEDs show that
the encoder works. Every LED is driven by a logic output (or comparator
output) through its own current-limiting resistor.\newpage}

\begin{description}
\item[{\rot\bf LED with series resistor}]~\\[-\lblskip]
\begin{center}
\begin{circuitikz}[scale=0.85, transform shape]
  \draw (0,2) node[left]{logic output $C_k$ or $B_i$ (\SI{5}{\volt})} to[short,o-] (0.5,2)
        to[R, l=$R$, i>^=$I_F$] (3,2) to[led, l=red LED] (3,0) node[ground]{};
  \draw (4.6,2) node[right]{} ;
\end{circuitikz}
\end{center}
\keyeq{Series resistor}{$\displaystyle R = \frac{V_\mathrm{OH} - V_F}{I_F} = \frac{\SI{5}{\volt}-\SI{2}{\volt}}{\SI{9}{\milli\ampere}} = \SI{333}{\ohm}\quad\to\quad\text{E12: }\SI{330}{\ohm},\ I_F = \SI{9.1}{\milli\ampere}$}
\bi
	\ite red LED: $V_F \approx$ \SIrange{1.8}{2.0}{\volt}; bright enough at \SIrange{5}{10}{\milli\ampere}
	\ite \textbf{one resistor per LED}: with a common resistor the current would depend on the number of LEDs that are on
	\ite Falstad LED model: $V_F = \SI{1.77}{\volt}$, $I_F = \SI{9.8}{\milli\ampere}$ with \SI{330}{\ohm}
\ei
\os{\newpage}

\item[{\rot\bf Current budget}]~\\[-\lblskip]
\bi
	\ite code 7: 7 bar LEDs + 3 binary LEDs on: $10\times\SI{9.1}{\milli\ampere} \approx \SI{91}{\milli\ampere}$ from the \SI{5}{\volt} supply
	\ite each output must deliver its LED current at a high level close to \SI{5}{\volt}
	\bi
		\itee in our circuit $C_4$ drives its bar LED, the $B_2$ LED and the gate inputs: \SI{18}{\milli\ampere}
		\itee 74HC gates: $\pm$\SI{25}{\milli\ampere} absolute maximum per output, $V_\mathrm{OH}$ drops under load
	\ei
	\ite better: buffer gates, low-current LEDs (\SI{2}{\milli\ampere}), or a driver IC (74HC245; ULN2003 sinks)
	\ite LED to ground (output sources current) or LED to \SI{5}{\volt} (output sinks, active low): both are fine
\ei
\nsl{The Falstad model has ideal outputs, so it does not show any of these
loading effects. In hardware, look up the guaranteed output voltage at the
required current in the data sheet. For a CMOS gate supplied with
\SI{5}{\volt}, $V_\mathrm{OH}$ is guaranteed above about \SI{3.8}{\volt}
only up to a few milliamperes; at higher currents the LED current
decreases accordingly.}
\os{\newpage}
\end{description}

\designcard{Falstad digital cheat-sheet}{%
\begin{itemize}\itemsep 0pt
\item gates: \emph{Draw $\to$ Logic Gates, Input and Output}: AND, OR, NAND, NOR, XOR, Inverter; number of inputs and high level (\SI{5}{\volt}) in the edit dialog (double click)
\item logic levels: input high above $V_\mathrm{high}/2 = \SI{2.5}{\volt}$; outputs are ideal \SI{0}{\volt}/\SI{5}{\volt} sources
\item comparator: op-amp with max output \SI{5}{\volt} and min output \SI{0}{\volt}; inputs $-$ and $+$ as drawn
\item \emph{Logic Input} (click to toggle) and \emph{Logic Output} (shows H/L) for quick tests
\item \emph{Labeled Node}: nodes with the same name are connected; use them for $C_1\ldots C_7$, $B_0\ldots B_2$; hover to read the voltage
\item input voltage: DC voltage source or rail with a slider (edit dialog); clock: rail with square wave, frequency, duty cycle
\item analog switch: closed if control $> \SI{2.5}{\volt}$, $R_\mathrm{on} = \SI{20}{\ohm}$
\item LED: always with a series resistor (\SI{330}{\ohm} at \SI{5}{\volt}); lights above a few mA
\item scope: right click on a node or element $\to$ \emph{View in Scope}; set the time step to \SI{0.1}{\milli\second} for \SI{5}{\hertz} clocks
\end{itemize}}
\os{\newpage}

%%%%%%%%%%%%%%%%%%%%%%%%%%%%%%%%%%%%%%%%%%%%%%%%%%%%%%%%%%%%%%%%%%%%%%%%%%%
\section{Other Converter Architectures}
%%%%%%%%%%%%%%%%%%%%%%%%%%%%%%%%%%%%%%%%%%%%%%%%%%%%%%%%%%%%%%%%%%%%%%%%%%%
\nsl{The flash converter is fast but expensive in comparators. Other
architectures trade speed for resolution. Most of them contain a
digital-to-analog converter (DAC) in a feedback loop, so we start with the
simplest DAC, the R-2R ladder. Then we look at the successive-approximation
converter found in every microcontroller, the dual-slope converter of
digital multimeters and the sigma-delta converter used for audio and
precision measurement \cite{kester2005,razavi1995}.\newpage}

\begin{description}
\item[{\rot\bf R-2R ladder DAC}]~\\[-\lblskip]
\begin{center}
\begin{circuitikz}[scale=0.8, transform shape]
  \draw (-1.8,3) to[R, l_=$2R$] (-1.8,0.8) node[ground]{};
  \draw (-1.8,3) -- (0,3) to[R, l=$R$] (2.5,3) to[R, l=$R$] (5,3) -- (6.8,3);
  \foreach \x/\b in {0/0,2.5/1,5/2} {
    \draw (\x,3) node[circ]{} to[R, l^=$2R$] (\x,0.8) node[below]{$B_\b$};
  }
  \draw (6.8,3) node[op amp, noinv input up, anchor=+](o){};
  \draw (o.-) -- ++(0,-0.8) -| (o.out);
  \draw (o.out) to[short,-o] ++(0.6,0) node[right]{$\Vout$};
\end{circuitikz}
\end{center}
\bi
	\ite bit inputs: logic levels \SI{0}{\volt} or $\Vref$; output resistance of the ladder: always $R$
	\ite $\displaystyle \Vout = \Vref\left(\frac{B_2}{2} + \frac{B_1}{4} + \frac{B_0}{8}\right) = \Vref\,\frac{D}{2^N}$
	\ite only two resistor values, any number of bits: easy to integrate
	\ite Falstad: $R = \SI{10}{\kilo\ohm}$, \SI{5}{\volt} logic, counter at \SI{1}{\hertz}/\SI{0.5}{\hertz}/\SI{0.25}{\hertz}: staircase in steps of \SI{0.625}{\volt}, $D = 5$: \SI{3.125}{\volt}
\ei
\falstad{r2r_dac}
\nsl{Every node of the ladder sees a resistance of $2R$ to the left (the
rest of the ladder) and $2R$ downwards (the bit input). Seen from
each bit input, the current therefore splits in half at every node on its
way to the output, so each bit has twice the weight of the next lower
bit. A Thevenin conversion from left to right gives the formula. Loading the output with
$4R$ instead of a follower scales it by $4R/(R+4R) = 0.8$; this is used in
the Falstad example \texttt{quant\_error} to get exactly
\SI{0.5}{\volt} per code from \SI{5}{\volt} logic levels.}
\os{\newpage}

\item[{\rot\bf Successive approximation (SAR)}]~\\[-\lblskip]
\begin{center}
\begin{tikzpicture}[font=\sffamily\small,>=Stealth,
	b/.style={draw,thick,minimum height=9mm,minimum width=18mm,align=center}]
	\node[b] (sh) at (0,0) {S\&H};
	\node[draw,thick,isosceles triangle,minimum width=12mm] (cmp) at (2.6,0) {};
	\node[b,fill=rwuviolet!12,minimum width=24mm] (sar) at (5.6,0) {SAR logic\\(register)};
	\node[b] (dac) at (2.6,-1.8) {DAC};
	\draw[->,thick] (-1.6,0) node[left]{$\Vin$} -- (sh);
	\draw[->,thick] (sh) -- (cmp.west |- 0,0.25) node[pos=0.8,above]{\scriptsize $+$};
	\draw[->,thick] (cmp) -- (sar);
	\draw[->,thick] (sar) |- (dac);
	\draw[->,thick] (dac.west) -- ++(-0.5,0) |- (cmp.west |- 0,-0.25);
	\draw[->,thick] (sar.east) -- ++(1,0) node[right]{$D$};
\end{tikzpicture}
\end{center}
\bi
	\ite binary search with one comparator and a DAC: one bit per clock cycle, MSB first
	\ite set the bit, compare $\Vin$ with $V_\mathrm{DAC}$: keep the bit if $\Vin \geq V_\mathrm{DAC}$, otherwise clear it
	\ite $N$ steps for $N$ bits; needs an S\&H (built in); 8 \ldots 18 bit, up to some MS/s
\ei
\os{\newpage}

\item[{\rot\bf SAR: worked example}]~\\[-\lblskip]
\begin{center}
\begin{tikzpicture}
\begin{axis}[width=6.4cm,height=4.4cm,xmin=0,xmax=4,ymin=0,ymax=4,
	xlabel={step},ylabel={$V$ / V},axis lines=left,font=\scriptsize,xtick={0,1,2,3,4},
	ytick={0,1,2,2.5,3,4}]
\addplot[rwuviolet,very thick,const plot] coordinates {(0,2) (1,3) (2,2.5) (3,2.75) (4,2.75)};
\addplot[rwucyan,thick,dashed] coordinates {(0,2.6) (4,2.6)};
\node[font=\tiny,anchor=south] at (axis cs:3.5,2.75) {$V_\mathrm{DAC}$};
\node[font=\tiny,anchor=north] at (axis cs:0.5,2.55) {$\Vin$};
\end{axis}
\end{tikzpicture}
\hspace{4mm}
\begin{tabular}[b]{cccc}
\toprule
step & trial & $V_\mathrm{DAC}$ & $\Vin \geq V_\mathrm{DAC}$?\\
\midrule
1 & \textbf{1}000 & \SI{2.00}{\volt} & yes: keep\\
2 & 1\textbf{1}00 & \SI{3.00}{\volt} & no: clear\\
3 & 10\textbf{1}0 & \SI{2.50}{\volt} & yes: keep\\
4 & 101\textbf{1} & \SI{2.75}{\volt} & no: clear\\
\midrule
 & \textbf{1010} & \SI{2.50}{\volt} & $D = 10$\\
\bottomrule
\end{tabular}
\end{center}
\bi
	\ite 4 bit, $V_\mathrm{FS} = \SI{4}{\volt}$, $\LSB = \SI{0.25}{\volt}$, $\Vin = \SI{2.6}{\volt}$
	\ite result 1010: $\Vin$ lies between \SI{2.50}{\volt} and \SI{2.75}{\volt}
\ei
\os{\newpage}

\item[{\rot\bf Dual-slope converter}]~\\[-\lblskip]
\begin{center}
\begin{circuitikz}[scale=0.72, transform shape]
  \draw (0,0) node[op amp](i){};
  \draw (i.+) -- ++(0,-0.4) node[ground]{};
  \draw (i.-) -- ++(-0.6,0) coordinate(m) to[R, l_=$R$] ++(-2,0) coordinate(sw);
  \draw (sw) -- ++(-0.4,0.6) node[left]{$\Vin$};
  \draw (sw) ++(-0.6,-0.5) node[left]{$-\Vref$};
  \draw (m) -- ++(0,1.4) coordinate(t) to[C, l=$C$] (t -| i.out) -- (i.out);
  \draw (i.out) -- ++(0.8,0) coordinate(o);
  \draw (o) ++(1.5,-0.5) node[op amp, noinv input up](k){};
  \draw (o) |- (k.+);
  \draw (k.-) -- ++(-0.3,0) node[ground, rotate=-90]{};
  \node[draw,thick,minimum width=2cm,minimum height=1cm,align=center] (cnt) at (6.6,-0.5) {control\\\& counter};
  \draw (k.out) -- (cnt.west |- k.out);
\end{circuitikz}
\begin{tikzpicture}
\begin{axis}[width=5.2cm,height=3.6cm,xmin=0,xmax=2.2,ymin=0,ymax=2.4,
	xlabel={$t$},ylabel={$-V_\mathrm{int}$},axis lines=left,font=\scriptsize,xtick={1},xticklabels={$T_1$},ytick=\empty]
\addplot[rwuviolet,thick] coordinates {(0,0) (1,2) (2,0)};
\addplot[rwucyan,thick] coordinates {(0,0) (1,1.2) (1.6,0)};
\end{axis}
\end{tikzpicture}
\end{center}
\bi
	\ite integrate $\Vin$ for a fixed time $T_1$, then integrate $-\Vref$ until the output is zero again ($T_2$)
	\ite $\Vin\,T_1 = \Vref\,T_2$ \quad$\Rightarrow$\quad $\displaystyle D = N_1\frac{\Vin}{\Vref}$: independent of $R$, $C$ and the clock frequency
	\ite $T_1 = \SI{20}{\milli\second}$ (or \SI{100}{\milli\second}): 50 Hz interference integrates to zero
\ei
\nsl{The dual-slope converter is slow (a few conversions per second) but
very accurate, and it rejects mains interference by design: the integral
of a \SI{50}{\hertz} sine over a whole number of periods is zero. That is
why digital multimeters use it. It is the opposite strategy to our
capstone, where the 50 Hz interference is removed by the anti-aliasing
filter before sampling.}
\os{\newpage}

\item[{\rot\bf Sigma-delta converter (concept)}]~\\[-\lblskip]
\begin{center}
\begin{tikzpicture}[font=\sffamily\small,>=Stealth,
	b/.style={draw,thick,minimum height=9mm,minimum width=16mm,align=center}]
	\node[draw,thick,circle,inner sep=1pt] (sum) at (0,0) {$\Sigma$};
	\node[b] (int) at (1.8,0) {$\int$};
	\node[b] (q) at (4,0) {1-bit\\comparator};
	\node[b] (dac) at (2.9,-1.6) {1-bit DAC};
	\node[b,fill=rwuviolet!12,minimum width=26mm] (dec) at (7.4,0) {digital filter\\\& decimation};
	\draw[->,thick] (-1.4,0) node[left]{$\Vin$} -- (sum);
	\draw[->,thick] (sum) -- (int);
	\draw[->,thick] (int) -- (q);
	\draw[->,thick] (q) -- node[below,pos=0.45]{\scriptsize bits} (dec);
	\draw[->,thick] (5.4,0) |- (dac);
	\draw[->,thick] (dac) -| node[pos=0.8,left]{$-$} (sum);
	\draw[->,thick] (dec) -- ++(2,0) node[right]{$D$};
	\node[font=\scriptsize] at (4,1) {clock $k\cdot f_s$};
\end{tikzpicture}
\end{center}
\bi
	\ite oversampling with a 1-bit quantizer in a feedback loop: the density of ones is proportional to $\Vin$
	\ite the loop shapes the quantization noise to high frequencies; a digital low-pass removes it
	\ite 16 \ldots 24 bit, audio and precision measurement (weighing scales, thermocouples)
	\ite the digital filter is also the anti-aliasing filter: the analog filter can be simple
\ei
\os{\newpage}

\item[{\rot\bf Comparison}]~\\[-\lblskip]
\begin{center}
\footnotesize
\begin{tabular}{@{}lllll@{}}
\toprule
 & flash & SAR & dual-slope & sigma-delta\\
\midrule
principle & $2^N-1$ comparators & binary search & integration & oversampling\\
steps per sample & 1 & $N$ & $2\cdot 2^N$ clocks & many (oversampling)\\
resolution & 4 \ldots 8 bit & 8 \ldots 18 bit & 12 \ldots 22 bit & 16 \ldots 24 bit\\
speed & up to GS/s & kS/s \ldots MS/s & 1 \ldots 100 S/s & Hz \ldots MHz bandwidth\\
cost, power & high (exponential) & low & low & medium\\
typical use & oscilloscopes, radio & microcontrollers & multimeters & audio, scales\\
\bottomrule
\end{tabular}
\end{center}
\bi
	\ite the capstone uses a 3-bit flash converter because every part of it is visible: ladder, comparators, logic
\ei
\os{\newpage}
\end{description}

%%%%%%%%%%%%%%%%%%%%%%%%%%%%%%%%%%%%%%%%%%%%%%%%%%%%%%%%%%%%%%%%%%%%%%%%%%%
\section{Exercises}
%%%%%%%%%%%%%%%%%%%%%%%%%%%%%%%%%%%%%%%%%%%%%%%%%%%%%%%%%%%%%%%%%%%%%%%%%%%

\exercise{Aliasing}\label{ex:adc-alias}
A system samples at $f_s = \SI{5}{\hertz}$.
\begin{talist}
\ta Calculate the alias frequencies of input components at \SI{2}{\hertz},
	\SI{4}{\hertz}, \SI{12}{\hertz}, \SI{49}{\hertz} and \SI{50}{\hertz}.
\ta The sensor signal is \SI{2}{\volt} DC with \SI{0.5}{\volt} of \SI{50}{\hertz} hum.
	Describe the sampled values for $f_s = \SI{5}{\hertz}$ and for $f_s = \SI{4.9}{\hertz}$.
\ta An RC low-pass with \SI{6.8}{\kilo\ohm} and \SI{10}{\micro\farad} is inserted before the S\&H stage.
	Calculate $f_c$ and the remaining hum amplitude. Compare with the LSB of
	the 3-bit converter (\SI{0.5}{\volt}).
\ta Check part b) in Falstad.
\end{talist}

\begin{solution}
\begin{talist}
\ta \SI{2}{\hertz}: below $f_s/2$, no alias. \SI{4}{\hertz}: $|4-5| = \SI{1}{\hertz}$.
	\SI{12}{\hertz}: $|12-2\cdot5| = \SI{2}{\hertz}$. \SI{49}{\hertz}: $|49-10\cdot5| = \SI{1}{\hertz}$.
	\SI{50}{\hertz}: $|50-10\cdot5| = \SI{0}{\hertz}$.
\ta $f_s = \SI{5}{\hertz}$: the hum is sampled at the same phase every time; all samples
	are equal, $\SI{2}{\volt} + \SI{0.5}{\volt}\sin\varphi_0$, i.e.\ a constant offset between
	\SI{-0.5}{\volt} and \SI{+0.5}{\volt}. $f_s = \SI{4.9}{\hertz}$: the phase advances by
	$0.204$ periods per sample; the samples follow a \SI{1}{\hertz} sine between \SI{1.5}{\volt}
	and \SI{2.5}{\volt}.
\ta $f_c = 1/(2\pi\cdot\SI{6.8}{\kilo\ohm}\cdot\SI{10}{\micro\farad}) = \SI{2.34}{\hertz}$.
	$|H(\SI{50}{\hertz})| = 1/\sqrt{1+(50/2.34)^2} = 0.047$ (\SI{-26.6}{\dB}): the hum is reduced to
	\SI{23}{\milli\volt}, less than $\frac{1}{20}\,\LSB$. The offset error disappears from the readings.
\ta \falstad{alias50}\newline \texttt{Vs\_5Hz} holds \SI{2.14}{\volt} at every sample (offset
	\SI{0.14}{\volt} for the phase in the example); \texttt{Vs\_4.9Hz} shows \SI{2.14}{\volt},
	\SI{2.50}{\volt}, \SI{2.14}{\volt}, \SI{1.58}{\volt}, \SI{1.62}{\volt}, \SI{2.20}{\volt}, \ldots
\end{talist}
\end{solution}

\exercise{Quantization}
A 3-bit converter has a full scale of \SI{4}{\volt} and thresholds at
$k\cdot\LSB$.
\begin{talist}
\ta Calculate the LSB and the codes for \SI{0.3}{\volt}, \SI{1.8}{\volt}, \SI{2.0}{\volt} and \SI{3.98}{\volt}.
\ta Which voltage range does the code 101 represent? What is the largest quantization error?
\ta Calculate the SNR of the ideal 3-bit converter. How many bits are needed for an SNR of at least \SI{60}{\dB}?
\ta Observe the quantization in Falstad.
\end{talist}

\begin{solution}
\begin{talist}
\ta $\LSB = \SI{0.5}{\volt}$. $D = \lfloor\Vin/\SI{0.5}{\volt}\rfloor$: \SI{0.3}{\volt} $\to$ 000,
	\SI{1.8}{\volt} $\to$ 011, \SI{2.0}{\volt} $\to$ 011 or 100 (exactly on the threshold, the
	comparator offset decides), \SI{3.98}{\volt} $\to$ 111.
\ta 101 = 5: \SIrange{2.5}{3.0}{\volt}. The error $\Vin - D\cdot\LSB$ is between 0 and
	\SI{0.5}{\volt} (1 LSB); relative to the bin centre \SI{2.75}{\volt} it is $\pm\SI{0.25}{\volt}$.
\ta $\mathrm{SNR} = 6.02\cdot 3 + 1.76 = \SI{19.8}{\dB}$.
	$6.02\,N + 1.76 \geq 60$ $\Rightarrow$ $N \geq 9.7$: 10 bits.
\ta \falstad{quant_error}\newline At $\Vin = \SI{1.8}{\volt}$ the DAC output is $V_q = \SI{1.5}{\volt}$,
	at \SI{3.98}{\volt} it is \SI{3.5}{\volt}; $\Vin - V_q$ stays between \SI{0}{\volt} and \SI{0.5}{\volt}.
\end{talist}
\end{solution}

\exercise{Designing the 3-bit flash converter}
Design the converter of the capstone: full scale \SI{4.00}{\volt}, 3 bit,
comparators $C_k = (\Vin > k\cdot\SI{0.5}{\volt})$.
\begin{talist}
\ta Dimension the reference ladder with \SI{1}{\kilo\ohm} resistors. Calculate the tap voltages,
	the ladder current and the power taken from the reference.
\ta The comparators have an input bias current of \SI{100}{\nano\ampere}. Estimate the
	worst-case threshold shift. (Hint: Thevenin resistance of a tap.)
\ta Write down the truth table from thermometer code to binary and derive the logic
	equations for $B_2$, $B_1$, $B_0$.
\ta How must the ladder be changed for thresholds at $(k-\frac12)\cdot\LSB$?
\ta Check the converter in Falstad at the bin centres.
\end{talist}

\begin{solution}
\begin{talist}
\ta 8 resistors of \SI{1}{\kilo\ohm} from \SI{4}{\volt} to ground. Taps
	$V_k = k\cdot\SI{4}{\volt}/8 = \SI{0.5}{\volt}, \SI{1.0}{\volt}, \ldots, \SI{3.5}{\volt}$.
	$I = \SI{4}{\volt}/\SI{8}{\kilo\ohm} = \SI{0.5}{\milli\ampere}$, $P = \SI{2}{\milli\watt}$.
\ta The largest Thevenin resistance is at the middle tap: $\SI{4}{\kilo\ohm}\parallelto\SI{4}{\kilo\ohm}
	= \SI{2}{\kilo\ohm}$. Shift: $\SI{100}{\nano\ampere}\cdot\SI{2}{\kilo\ohm} = \SI{0.2}{\milli\volt}$,
	negligible compared with \SI{0.5}{\volt}.
\ta Truth table as on the slide \emph{Thermometer code}. With $H_k = C_k\overline{C_{k+1}}$:
	$B_2 = H_4+H_5+H_6+H_7 = C_4$; $B_1 = H_2+H_3+H_6+H_7 = C_2\overline{C_4} + C_6$;
	$B_0 = H_1+H_3+H_5+H_7 = C_1\overline{C_2} + C_3\overline{C_4} + C_5\overline{C_6} + C_7$.
\ta Total resistance stays $8R$: bottom resistor $R/2 = \SI{500}{\ohm}$, six resistors $R$, top
	resistor $3R/2 = \SI{1.5}{\kilo\ohm}$. The taps are then \SI{0.25}{\volt}, \SI{0.75}{\volt}, \ldots, \SI{3.25}{\volt}.
\ta \falstad{flash3}\newline Taps \SI{0.5}{\volt} \ldots \SI{3.5}{\volt}, ladder current
	\SI{0.50}{\milli\ampere}. Slider at \SI{0.25}{\volt}, \SI{0.75}{\volt}, \ldots, \SI{3.75}{\volt}:
	the bar graph shows 0 \ldots 7 LEDs, the binary LEDs 000, 001, \ldots, 111.
\end{talist}
\end{solution}

\exercise{Bubble error}
Because of noise, the comparators of the 3-bit flash converter deliver the
pattern 0010111 ($C_7$ first), a bubble at $C_4$.
\begin{talist}
\ta Which input range is plausible? Which code would be correct?
\ta Which code does the encoder of the capstone produce?
\ta Which code does a priority encoder produce?
\end{talist}

\begin{solution}
\begin{talist}
\ta $C_5 = 1$ but $C_4 = 0$ contradict each other. Either $C_4$ is wrong (input
	between \SI{2.5}{\volt} and \SI{3.0}{\volt}, code 5) or $C_5$ is wrong (input between
	\SI{1.5}{\volt} and \SI{2.0}{\volt}, code 3).
\ta $B_2 = C_4 = 0$; $B_1 = C_2\overline{C_4} + C_6 = 1$; $B_0 = C_1\overline{C_2} + C_3\overline{C_4} +
	C_5\overline{C_6} + C_7 = 0+1+1+0 = 1$: code 011 = 3, one of the plausible values. (For the pattern 0000101 of the slide,
	however, it gives 001 instead of 011.)
\ta The priority encoder looks at the highest active input, $C_5$: code 101 = 5, the other
	plausible value. For 0000101 it gives 011 = 3, the correct value; it never produces a code
	below the top of the column.
\end{talist}
\end{solution}

\exercise{LED display}
The bar graph and the binary LEDs are driven by \SI{5}{\volt} logic outputs.
Red LEDs with $V_F \approx \SI{2}{\volt}$ shall carry \SIrange{5}{10}{\milli\ampere}.
\begin{talist}
\ta Choose the series resistor from the E12 series and calculate the LED current.
\ta Calculate the total LED current for the code 111.
\ta In the Falstad model, $C_4$ drives its bar LED and also the $B_2$ LED. Which current does
	this output deliver?
\ta Why must each LED have its own resistor?
\end{talist}

\begin{solution}
\begin{talist}
\ta $R = (\SI{5}{\volt}-\SI{2}{\volt})/I_F$ gives \SI{300}{\ohm} \ldots \SI{600}{\ohm}. Choose \SI{330}{\ohm}:
	$I_F = \SI{3}{\volt}/\SI{330}{\ohm} = \SI{9.1}{\milli\ampere}$.\newline
	\falstad{flash3}\newline The Falstad LED has $V_F = \SI{1.77}{\volt}$, so the current is
	$\SI{3.23}{\volt}/\SI{330}{\ohm} = \SI{9.8}{\milli\ampere}$.
\ta 10 LEDs: $10\cdot\SI{9.1}{\milli\ampere} = \SI{91}{\milli\ampere}$.
\ta $2\cdot\SI{9.1}{\milli\ampere} = \SI{18}{\milli\ampere}$ (plus negligible gate input currents).
	That is near the limit of a CMOS logic output; a buffer or lower LED currents
	(\SI{680}{\ohm}: \SI{4.4}{\milli\ampere}) are safer.
\ta LEDs of different colours and samples have different $V_F$. With a common resistor the
	LED with the lowest $V_F$ takes most of the current, and the current per LED changes with
	the number of LEDs that are on.
\end{talist}
\end{solution}

\exercise{Schmitt trigger, R-2R DAC and SAR}
\begin{talist}
\ta A non-inverting Schmitt trigger has $\Vref = \SI{2}{\volt}$, $R_1 = \SI{10}{\kilo\ohm}$ and
	output levels \SI{0}{\volt}/\SI{5}{\volt}. Choose $R_2$ for a hysteresis of \SI{0.5}{\volt} and
	calculate both thresholds.
\ta A 3-bit R-2R DAC is driven by \SI{5}{\volt} logic levels. Calculate the step size and the output
	for $D = 5$.
\ta A 4-bit SAR converter ($V_\mathrm{FS} = \SI{4}{\volt}$) converts $\Vin = \SI{1.3}{\volt}$.
	Write down the trial codes, the DAC voltages and the result.
\ta Check a) and b) in Falstad.
\end{talist}

\begin{solution}
\begin{talist}
\ta Hysteresis $V_\mathrm{OH}R_1/R_2 = \SI{0.5}{\volt}$ $\Rightarrow$ $R_2 = \SI{5}{\volt}\cdot\SI{10}{\kilo\ohm}/\SI{0.5}{\volt}
	= \SI{100}{\kilo\ohm}$. $V_{T+} = \SI{2}{\volt}\cdot 1.1 = \SI{2.2}{\volt}$, $V_{T-} = \SI{1.7}{\volt}$.
\ta $\LSB = \SI{5}{\volt}/8 = \SI{0.625}{\volt}$; $D = 5$: $\SI{3.125}{\volt}$.
\ta $\LSB = \SI{0.25}{\volt}$. 1000: \SI{2.00}{\volt}, no $\to$ 0. 0100: \SI{1.00}{\volt}, yes $\to$ 1.
	0110: \SI{1.50}{\volt}, no $\to$ 0. 0101: \SI{1.25}{\volt}, yes $\to$ 1. Result 0101 = 5
	(\SIrange{1.25}{1.50}{\volt}).
\ta \falstad{comparator_hyst}\newline The Schmitt trigger switches at \SI{2.2}{\volt} (rising) and
	\SI{1.7}{\volt} (falling) and gives exactly two edges per period; the plain comparator
	produces 10 edges.\newline
	\falstad{r2r_dac}\newline \texttt{Vout} steps through \SI{0}{\volt}, \SI{0.625}{\volt}, \ldots,
	\SI{4.375}{\volt}; at $D = 5$ it shows \SI{3.125}{\volt}.
\end{talist}
\end{solution}

%%% Local Variables:
%%% mode: latex
%%% TeX-master: "docu"
%%% End:
